\section*{अध्याय \ref{ch:g5:areas} --- परिमाप और क्षेत्रफल}

\begin{solution}{exo:g5:areas:1}
\omterm{def:g4:measure:area}{क्षेत्रफल}: $8 \times 5 = 40$ cm$^2$। \omterm{def:g4:measure:perimeter}{परिमाप}:
$2 \times (8 + 5) = 26$ cm।
\end{solution}

\begin{solution}{exo:g5:areas:2}
$9 \times 9 = 81$ cm$^2$; \quad $12 \times 7 = 84$ m$^2$; \quad
$4.5 \times 6 = 27$ cm$^2$।
\end{solution}

\begin{solution}{exo:g5:areas:3}
चौड़ाई: $63 \div 9 = 7$ cm। \omterm{def:g4:measure:perimeter}{परिमाप}: $2 \times (9 + 7) = 32$ cm।
\end{solution}

\begin{solution}{exo:g5:areas:4}
जैसे $6 \times 4$ (\omterm{def:g4:measure:perimeter}{परिमाप} $20$) और $8 \times 3$
(\omterm{def:g4:measure:perimeter}{परिमाप} $22$): \omterm{def:g4:measure:area}{क्षेत्रफल} वही $24$, \omterm{def:g4:measure:perimeter}{परिमाप} अलग।
\end{solution}

\begin{solution}{exo:g5:areas:5}
$3$ m$^2 = 30\,000$ cm$^2$; \quad $50\,000$ cm$^2 = 5$ m$^2$।
\end{solution}

\begin{solution}{exo:g5:areas:6}
आड़ी पट्टी: $8 \times 2 = 16$ cm$^2$; खड़ी: $2 \times 5 = 10$ cm$^2$;
कुल $26$ cm$^2$।
\end{solution}

\begin{solution}{exo:g5:areas:7}
$30 \times 20 - 24 \times 14 = 600 - 336 = 264$ cm$^2$ लकड़ी।
\end{solution}

\begin{solution}{exo:g5:areas:8}
$8 \times 6$ के आयत का \omterm{def:g3:division:half}{आधा}: $48 \div 2 = 24$ cm$^2$।
\end{solution}

\begin{solution}{exo:g5:areas:9}
बीज: \omterm{def:g4:measure:area}{क्षेत्रफल} $7 \times 4 = 28$ m$^2$, ख़र्च $28 \times 2 = 56$।
बाड़: \omterm{def:g4:measure:perimeter}{परिमाप} $2 \times (7 + 4) = 22$ m, ख़र्च
$22 \times 3 = 66$।
\end{solution}

\begin{solution}{exo:g5:areas:10}
$50$ cm की दो टाइलें मिलकर एक मीटर बनाती हैं: $12$ m लंबाई पर
$24$ टाइलें; $2$ m चौड़ाई पर $4$ टाइलें। कुल:
$24 \times 4 = 96$ टाइलें।
\end{solution}

\begin{solution}{exo:g5:areas:11}
पहले: \omterm{def:g4:measure:perimeter}{परिमाप} $2 \times (4 + 3) = 14$ cm, \omterm{def:g4:measure:area}{क्षेत्रफल} $12$ cm$^2$।
बाद में (भुजाएँ $8$ और $6$): \omterm{def:g4:measure:perimeter}{परिमाप} $28$ cm, \omterm{def:g4:measure:area}{क्षेत्रफल} $48$ cm$^2$।
\omterm{def:g4:measure:perimeter}{परिमाप} \emph{दुगुना} हुआ; \omterm{def:g4:measure:area}{क्षेत्रफल} \emph{चार} गुना
($2 \times 2$): लंबाइयाँ और \omterm{def:g4:measure:area}{क्षेत्रफल} एक ही ढंग से नहीं
बढ़ते।
\end{solution}
